\subsection{DINOv3 ViT-B/16 Fine-Tuned: Full Results}
\label{app:dinov3_full}

Like the two detectors above, the fine-tuned DINOv3 ViT-B/16 results of Tab.~\ref{tab:dinov3_main} come with a cross-domain transfer table (Tab.~\ref{tab:dinov3_ft_cd}) and a per-tier STRANGE-Bench breakdown (Tab.~\ref{tab:dinov3_ft_strange}). The splits are again the quintile ones, and every number is F1\textsubscript{micro} (the natural metric for the multi-label classification head; see \S\ref{app:exp-recipes}). The augmentation suite mirrors that of the detection families.

\noindent\textbf{In-domain typicality.}
On COCO (Tab.~\ref{tab:dinov3_ft_cd}, left block), Std.\ aug yields a typicality gap of $\Delta^{\text{T-A}} = {}$$5.5$ at $\text{F1}^{\text{val}} = {}$$81.9$. Color jitter and MixUp shave the gap modestly (4.6 and 4.2) without hurting overall F1. Mosaic and YOLOX produce the smallest non-paste gaps ($4.0$ and $3.7$) but \emph{collapse} overall F1 to $69.6$ and $70.0$, the same uniform-degradation trap visible in the detection families (\S\ref{app:codino_full}). XPaste\,+\,AAS \emph{(ours)} attains the highest in-domain F1 ($82.5$) and a gap of $\Delta^{\text{T-A}} = 3.9$ ($1.6$ below Std.; SCP reaches $3.7$).

\noindent\textbf{Cross-domain typicality.}
On Objects365, for the right block (pairs ranked by NPMI recomputed on Objects365 train, \S\ref{app:exp-datasets}, i.e.\ by the co-occurrence statistics of the domain the images come from), the gap is $7.1$ to $8.0$~p.p.\@ for every row except Mosaic and YOLOX, against AP gaps of $3.8$ to $5.1$ for Co-DINO (Tab.~\ref{tab:codino_cd}) and $3.0$ to $3.5$ for GDINO (Tab.~\ref{tab:gdino_o365}) on the same split, so the typicality bias of the fine-tuned classifier carries over to real images of another dataset. Mosaic and YOLOX have smaller gaps ($4.8$ and $4.6$) only because both arms collapse (typical F1 $26.8$ and $27.3$ against $38.1$ to $41.2$ for the other rows), as on COCO. SCP and SCP\,+\,AAS \emph{(ours)} have the highest typical and atypical F1 ($41.2 / 33.7$ and $41.0 / 33.6$); against Std.\ ($40.1 / 32.2$) both raise the atypical arm more than the typical one ($+1.5$ vs.\ $+1.1$ and $+1.4$ vs.\ $+0.9$), so the gap narrows from $7.9$ to $7.5$ and $7.4$, the smallest apart from Mosaic, YOLOX and No aug ($7.1$, at lower F1 on both arms). XPaste and XPaste\,+\,AAS stay at the Std.\ level on both arms ($40.0 / 32.2$ and $40.0 / 32.3$), so the in-domain lead of XPaste\,+\,AAS does not carry into this transfer. With COCO NPMI (middle block), every Objects365 gap is $1.7$ to $1.8$ times larger ($12.6$ to $14.2$~p.p.\@; Mosaic and YOLOX $8.0$ and $8.1$), whereas the Co-DINO and GDINO gaps are smaller there ($2.0$ to $3.0$ and $2.1$ to $2.7$); SCP and SCP\,+\,AAS again have the highest atypical F1 ($31.4$ and $31.1$).

\begin{table*}[!htp]%
\centering
\footnotesize
\setlength{\tabcolsep}{4pt}
\caption{\textbf{DINOv3 ViT-B/16 fine-tuned: in-domain (COCO) and
cross-domain (Objects365) typicality.} On COCO, XPaste\,+\,AAS \emph{(ours)} has the highest F1$^{\text{val}}$, typical and atypical F1. Ranked by Objects365's own NPMI (right block), Objects365 shows a gap of $7.1$ to $8.0$~p.p.\@ for every row except Mosaic and YOLOX, whose arms collapse, and SCP and SCP\,+\,AAS \emph{(ours)} have the highest typical and atypical F1, raising the atypical arm more than the typical one over Std. Ranked by COCO NPMI (middle block), every Objects365 gap is $1.7$ to $1.8$ times larger. F1$_{\text{micro}}$ (\%) on the quintile paired splits of COCO val and Objects365 val, the latter split by COCO \texttt{train2017} NPMI or by NPMI recomputed on Objects365 train (App.~\ref{app:exp-datasets}); rows as in Tab.~\ref{tab:dinov3_main} without the linear probe. Best per column in \textbf{bold}, second-best \underline{underlined}.}
\label{tab:dinov3_ft_cd}
\begin{tabular}{l cccc ccc ccc}
\toprule
& \multicolumn{4}{c}{\textbf{COCO}} & \multicolumn{3}{c}{\textbf{Objects365, COCO NPMI}} & \multicolumn{3}{c}{\textbf{Objects365, O365 NPMI}} \\
\cmidrule(lr){2-5} \cmidrule(lr){6-8} \cmidrule(lr){9-11}
Augmentation & $\text{F1}^{\text{val}}$ $\uparrow$ & $\text{F1}^{\text{T}}$ $\uparrow$ & $\text{F1}^{\text{A}}$ $\uparrow$ & $\Delta^{\text{T-A}}$ $\downarrow$ & $\text{F1}^{\text{T}}$ $\uparrow$ & $\text{F1}^{\text{A}}$ $\uparrow$ & $\Delta^{\text{T-A}}$ $\downarrow$ & $\text{F1}^{\text{T}}$ $\uparrow$ & $\text{F1}^{\text{A}}$ $\uparrow$ & $\Delta^{\text{T-A}}$ $\downarrow$ \\
\midrule
No aug & 78.9 & 77.8 & 70.1 & 7.7 & 42.0 & 29.3 & 12.6 & 38.1 & 30.9 & 7.1 \\
\midrule
\multicolumn{11}{l}{\emph{Non-generative}} \\
Std. & 81.9 & 79.8 & 74.3 & 5.5 & 43.9 & 30.1 & 13.8 & 40.1 & 32.2 & 7.9 \\
Color jitter & 81.8 & 79.7 & 75.1 & 4.6 & 43.8 & 30.1 & 13.7 & 40.2 & 32.2 & 8.0 \\
MixUp & \underline{82.2} & 79.6 & 75.4 & 4.2 & 43.9 & 30.2 & 13.7 & 40.1 & 32.6 & 7.5 \\
Mosaic & 69.6 & 71.0 & 67.0 & 4.0 & 28.6 & 20.5 & \textbf{8.0} & 26.8 & 22.0 & \underline{4.8} \\
YOLOX & 70.0 & 71.2 & 67.4 & \textbf{3.7} & 28.9 & 20.8 & \underline{8.1} & 27.3 & 22.7 & \textbf{4.6} \\
SCP & 81.6 & 78.8 & 75.1 & \textbf{3.7} & \textbf{44.2} & \textbf{31.4} & 12.8 & \textbf{41.2} & \textbf{33.7} & 7.5 \\
\midrule
\multicolumn{11}{l}{\emph{Generative}} \\
InstAug & 81.9 & 79.5 & 74.4 & 5.2 & \underline{44.0} & 30.1 & 13.9 & 40.3 & 32.5 & 7.9 \\
XPaste & 82.1 & \underline{80.2} & \underline{76.1} & 4.1 & 43.8 & 29.7 & 14.2 & 40.0 & 32.2 & 7.8 \\
\midrule
\rowcolor{oursrow} SCP + AAS \emph{(ours)} & 81.7 & 79.0 & 75.0 & 4.0 & \underline{44.0} & \underline{31.1} & 12.8 & \underline{41.0} & \underline{33.6} & 7.4 \\
\rowcolor{oursrow} \textbf{XPaste + AAS \emph{(ours)}} & \textbf{82.5} & \textbf{80.6} & \textbf{76.8} & \underline{3.9} & 43.6 & 29.7 & 13.9 & 40.0 & 32.3 & 7.7 \\
\bottomrule
\end{tabular}

\end{table*}

\noindent\textbf{Per-tier STRANGE-Bench breakdown.} Tab.~\ref{tab:dinov3_ft_strange} reports the full per-tier STRANGE-Bench breakdown for DINOv3 ViT-B/16 FT. Three observations stand out. First, this family has the largest single-tier gap in the paper: the Std.\ baseline gap grows from $\Delta = -0.7$ at $T_1$ to $\Delta = {}$$40.6$ at $T_5$, i.e.\ a $41$ p.p.\@ swing as contextual load increases, with $T_{\text{never}}$ at only $62.6$. The mean gap $\Delta_{1\text{-}5} = {}$$22.2$ is roughly $3.5\times$ that of Co-DINO Std.\ ($6.3$) and about $9\times$ that of GDINO zero-shot ($2.4$), reflecting the strong typicality prior baked into a multi-label head trained on natural co-occurrences. Second, the paste-based interventions are correspondingly more dramatic: SCP alone reduces the mean gap to $4.3$ and lifts $T_{\text{never}}$ to $86.5$ (a $+24$ p.p.\@ absolute improvement over Std.). XPaste\,+\,AAS \emph{(ours)} pushes both numbers further to $\Delta_{1\text{-}5} = {}$$-0.8$ and $T_{\text{never}} = {}$$96.1$ (XPaste alone: $-0.7$ and $95.0$), the largest mitigation effect we observe in this paper. Third, the gap-closure profile is itself revealing: under XPaste\,+\,AAS the per-tier gap stays at or below $0.7$~p.p.\@ for every tier $T_1$ to $T_5$ (negative at $T_1$ to $T_3$, where atypical F1 is higher), including $T_5$ (gap $= {}$$0.7$, vs $40.6$ for Std.\ and $33.8$ for No-aug). This means the model's reliance on typical co-occurrence in the highest-load regime is essentially eliminated, while typical-tier F1 stays well above Std. ($\text{T}5^{\text{T}} = {}$$93.3$ vs Std.'s $88.6$). Mosaic and YOLOX, despite small mean gaps ($6.7$ and $6.0$), again pay for it with collapsed typical scores ($T^{\text{T}}_{1\text{-}5} = 74.7$ and $75.1$ vs.\ $89.7$ for Std.).

\begin{table*}[!htp]%
\centering
\scriptsize
\setlength{\tabcolsep}{2pt}
\caption{\textbf{DINOv3 ViT-B/16 fine-tuned: STRANGE-Bench per-tier breakdown.} The Std.\ gap grows from $-0.7$ at T1 to $40.6$ at T5 (mean $22.2$), the largest single-tier gap in the paper; SCP brings the mean down to $4.3$ and XPaste\,+\,AAS \emph{(ours)} to $-0.8$, with gaps between $-2.4$ and $0.7$ at every tier and the highest $T_{\text{never}}$ ($96.1$). F1$_{\text{micro}}$ (\%) of the typical (T) and atypical (A) arm of each tier and their gap $\Delta$; $T_{1\text{-}5}$ is the mean over T1 to T5, rows as in Tab.~\ref{tab:dinov3_main} without the linear probe. Best per column in \textbf{bold}, second-best \underline{underlined}.}
\label{tab:dinov3_ft_strange}
\resizebox{\linewidth}{!}{%
\begin{tabular}{l ccc ccc ccc ccc ccc ccc c}
\toprule
& \multicolumn{3}{c}{\textbf{Tier 1}} & \multicolumn{3}{c}{\textbf{Tier 2}} & \multicolumn{3}{c}{\textbf{Tier 3}} & \multicolumn{3}{c}{\textbf{Tier 4}} & \multicolumn{3}{c}{\textbf{Tier 5}} & \multicolumn{3}{c}{\textbf{Mean T1--T5}} & \\
\cmidrule(lr){2-4} \cmidrule(lr){5-7} \cmidrule(lr){8-10} \cmidrule(lr){11-13} \cmidrule(lr){14-16} \cmidrule(lr){17-19}
Aug & T$\uparrow$ & A$\uparrow$ & $\Delta\downarrow$ & T$\uparrow$ & A$\uparrow$ & $\Delta\downarrow$ & T$\uparrow$ & A$\uparrow$ & $\Delta\downarrow$ & T$\uparrow$ & A$\uparrow$ & $\Delta\downarrow$ & T$\uparrow$ & A$\uparrow$ & $\Delta\downarrow$ & T$\uparrow$ & A$\uparrow$ & $\Delta\downarrow$ & $T_{\text{never}}\uparrow$ \\
\midrule
No aug & 87.5 & 89.2 & $-1.7$ & 87.2 & 75.6 & 11.6 & 85.5 & 63.2 & 22.3 & 84.4 & 54.4 & 30.1 & 82.8 & 49.1 & 33.8 & 85.5 & 66.3 & 19.2 & 61.5 \\
\midrule
\multicolumn{20}{l}{\emph{Non-generative}} \\
Std. & 91.3 & 92.0 & $-0.7$ & 91.0 & 79.2 & 11.7 & 89.2 & 64.2 & 25.0 & 88.3 & 54.0 & 34.3 & 88.6 & 48.0 & 40.6 & 89.7 & 67.5 & 22.2 & 62.6 \\
Color jitter & 91.6 & 91.5 & 0.1 & 90.6 & 79.6 & 11.0 & 89.2 & 64.3 & 24.8 & 88.3 & 52.4 & 35.9 & 86.5 & 46.0 & 40.6 & 89.2 & 66.7 & 22.5 & 62.7 \\
MixUp & 91.4 & 93.5 & $-2.2$ & 91.0 & 85.6 & 5.4 & 90.2 & 70.8 & 19.4 & 89.1 & 60.5 & 28.5 & 88.0 & 51.6 & 36.4 & 89.9 & 72.4 & 17.5 & 69.9 \\
Mosaic & 57.8 & 51.5 & 6.3 & 71.7 & 70.2 & 1.5 & 76.9 & 73.3 & 3.6 & 83.1 & 73.0 & 10.1 & 84.2 & 72.1 & 12.1 & 74.7 & 68.0 & 6.7 & 70.6 \\
YOLOX & 59.2 & 55.0 & 4.2 & 71.9 & 71.7 & 0.1 & 78.2 & 73.4 & 4.9 & 82.4 & 73.3 & 9.1 & 84.0 & 72.1 & 11.9 & 75.1 & 69.1 & 6.0 & 70.7 \\
SCP & 92.3 & 94.0 & $-1.7$ & 92.4 & 92.5 & $-0.1$ & 91.5 & 89.0 & 2.5 & 91.6 & 83.4 & 8.1 & 89.4 & 76.7 & 12.7 & 91.4 & 87.1 & 4.3 & 86.5 \\
\midrule
\multicolumn{20}{l}{\emph{Generative}} \\
InstAug & 91.4 & 93.6 & $-2.2$ & 91.3 & 79.1 & 12.2 & 89.1 & 64.2 & 24.9 & 88.6 & 53.1 & 35.5 & 87.2 & 48.4 & 38.7 & 89.5 & 67.7 & 21.8 & 63.3 \\
XPaste & 93.5 & \underline{96.7} & \textbf{\boldmath $-3.1$} & \underline{93.8} & \underline{96.4} & \textbf{\boldmath $-2.6$} & \underline{93.9} & \underline{94.0} & \underline{$-0.1$} & \underline{94.1} & \underline{93.6} & \underline{0.5} & \textbf{93.4} & \underline{91.4} & \underline{1.9} & \underline{93.8} & \underline{94.4} & \underline{$-0.7$} & \underline{95.0} \\
\midrule
\rowcolor{oursrow} SCP + AAS \emph{(ours)} & \underline{94.2} & 96.3 & $-2.1$ & 92.9 & 94.7 & \underline{$-1.8$} & 92.2 & 91.6 & 0.6 & 92.3 & 88.8 & 3.5 & 90.8 & 84.8 & 6.0 & 92.5 & 91.2 & 1.2 & 89.5 \\
\rowcolor{oursrow} \textbf{XPaste + AAS \emph{(ours)}} & \textbf{94.9} & \textbf{97.3} & \underline{$-2.4$} & \textbf{95.4} & \textbf{97.2} & \underline{$-1.8$} & \textbf{94.7} & \textbf{95.5} & \textbf{\boldmath $-0.9$} & \textbf{94.9} & \textbf{94.6} & \textbf{0.3} & \underline{93.3} & \textbf{92.6} & \textbf{0.7} & \textbf{94.6} & \textbf{95.5} & \textbf{\boldmath $-0.8$} & \textbf{96.1} \\
\bottomrule
\end{tabular}
}
\end{table*}
